% from ega2-indexes-and-contents-fr.tex
\oldpage[II]{214}
\markboth{A. GROTHENDIECK --- Chap. II}{A. GROTHENDIECK --- Chap. II}
\begin{longtable}{@{}p{0.09\textwidth}p{0.69\textwidth}r@{}}
\hyperref[section:II.4-fr]{\textsection4} & Fibrés projectifs. Faisceaux amples & 71\\
\hyperref[subsection:II.4.1-fr]{4.1} & Définition des fibrés projectifs & 71\\
\hyperref[subsection:II.4.2-fr]{4.2} & Morphismes d'un préschéma dans un fibré projectif & 72\\
\hyperref[subsection:II.4.3-fr]{4.3} & Le morphisme de Segre & 76\\
\hyperref[subsection:II.4.4-fr]{4.4} & Immersions dans les fibrés projectifs. Faisceaux très amples & 78\\
\hyperref[subsection:II.4.5-fr]{4.5} & Faisceaux amples & 83\\
\hyperref[II.4.6-fr]{4.6} & Faisceaux relativement amples & 89\\
\hyperref[section:II.5-fr]{\textsection5} & Morphismes quasi-affines; morphismes quasi-projectifs; morphismes propres; morphismes projectifs & 94\\
\hyperref[subsection:II.5.1-fr]{5.1} & Morphismes quasi-affines & 94\\
\hyperref[subsection:II.5.2-fr]{5.2} & Le critère de Serre & 97\\
\hyperref[subsection:II.5.3-fr]{5.3} & Morphismes quasi-projectifs & 99\\
\hyperref[subsection:II.5.4-fr]{5.4} & Morphismes propres et morphismes universellement fermés & 100\\
\hyperref[subsection:II.5.5-fr]{5.5} & Morphismes projectifs & 103\\
\hyperref[subsection:II.5.6-fr]{5.6} & Le lemme de Chow & 106\\
\hyperref[section:II.6-fr]{\textsection6} & Morphismes entiers et morphismes finis & 110\\
\hyperref[subsection:II.6.1-fr]{6.1} & Préschémas entiers sur un autre & 110\\
\hyperref[subsection:II.6.2-fr]{6.2} & Morphismes quasi-finis & 114\\
\hyperref[subsection:II.6.3-fr]{6.3} & Fermeture intégrale d'un préschéma & 116\\
\hyperref[subsection:II.6.4-fr]{6.4} & Déterminant d'un endomorphisme de $\mathcal{O}_X$-Module & 120\\
\hyperref[subsection:II.6.5-fr]{6.5} & Norme d'un faisceau inversible & 125\\
\hyperref[subsection:II.6.6-fr]{6.6} & Application : critères d'amplitude & 130\\
\hyperref[subsection:II.6.7-fr]{6.7} & Le théorème de Chevalley & 135\\
\hyperref[section:II.7-fr]{\textsection7} & Critères valuatifs & 138\\
\hyperref[subsection:II.7.1-fr]{7.1} & Rappels sur les anneaux de valuation & 138\\
\hyperref[subsection:II.7.2-fr]{7.2} & Critère valuatif de séparation & 141\\
\hyperref[subsection:II.7.3-fr]{7.3} & Critère valuatif de propreté & 143\\
\hyperref[subsection:II.7.4-fr]{7.4} & Courbes algébriques et corps de fonctions de dimension 1 & 148\\
\hyperref[section:II.8-fr]{\textsection8} & Schémas éclatés; cônes projetants; fermeture projective & 152\\
\hyperref[subsection:II.8.1-fr]{8.1} & Préschémas éclatés & 152\\
\hyperref[subsection:II.8.2-fr]{8.2} & Résultats préliminaires sur la localisation dans les anneaux gradués & 157\\
\hyperref[subsection:II.8.3-fr]{8.3} & Cônes projetants & 162\\
\hyperref[subsection:II.8.4-fr]{8.4} & Fermeture projective d'un fibré vectoriel & 168\\
\hyperref[subsection:II.8.5-fr]{8.5} & Comportements fonctoriels & 169\\
\hyperref[subsection:II.8.6-fr]{8.6} & Un isomorphisme canonique pour les cônes épointés & 171\\
\end{longtable}
